\باب{جزئیتی تقسیم کا قانون برائے سمتی ضرب   }\شناخت{ضمیمہ_ز}

اس ضمیمہ میں    درج ذیل جزئیتی تقسیم کا قانون  (حصہ \حوالہء{10.4}، مساوات \حوالہء{6})  ثابت کیا جائے گا۔
\begin{align}\label{مساوات_ضمیمہ_جزئیتی_تقسیم_الف}
\kvec{A}\times (\kvec{B}+\kvec{C})=\kvec{A}\times \kvec{B}+\kvec{A}\times\kvec{C}
\end{align}

\ابتدا{ثبوتءء}\موٹا{ثبوت}\\
مساوات \حوالہ{مساوات_ضمیمہ_جزئیتی_تقسیم_الف} ثابت کرنے کی خاطر ہم \عددی{\kvec{A}\times\kvec{B}} کا خاکہ نئے انداز سے  بناتے ہیں۔ ہم مشترک نقطہ \عددی{O} سے \عددی{\kvec{A}} اور \عددی{\kvec{B}} کھینچ کر مستوی \عددی{M}  ،نقطہ   \عددی{O}  پر \عددی{\kvec{A}} کو قائمہ ، بناتے ہیں (شکل \حوالہء{A11})۔  ہم \عددی{M} پر \عددی{\kvec{B}} کا قائمہ تظلیل لیتے ہیں جو سمتیہ \عددی{\kvec{B}'}  ہو گا،   اور جس کی لمبائی  \عددی{|\kvec{B}|\sin\theta} ہو گی۔ ہم \عددی{\kvec{B}'}  کو \عددی{\kvec{A}} کے گرد مثبت سمت میں  \عددی{90^{\circ}} گھما کر سمتیہ \عددی{\kvec{B}''} حاصل کرتے ہیں۔آخر میں ہم   \عددی{\kvec{B}''} کو \عددی{\kvec{A}} کی لمبائی سے ضرب دیتے ہیں۔یوں حاصل سمتیہ \عددی{|\kvec{A}|\kvec{B}''}   درحقیقت \عددی{\kvec{A}\times\kvec{B}} کے برابر ہو گا، چونکہ اس کی بناؤٹ (شکل \حوالہء{A11})   یوں کی گئی کہ  \عددی{\kvec{B}''} کا رخ  وہی ہو جو  \عددی{\kvec{A}\times\kvec{B}}  کا رخ  ہے اور ساتھ ہی درج ذیل بھی ہے۔
\[|\kvec{A}||\kvec{B}''|=|\kvec{A}|\kvec{B}'|=|\kvec{A}||\kvec{B}|\sin\theta=|\kvec{A}\times\kvec{B}|\]
اب درج  ذیل میں سے ہر ایک عمل  ایسے   مثلث پر کرنے سے،   جس کا مستوی \عددی{\kvec{A}} کو عمودی نہ ہو،  ایک نیا مثلث پیدا ہو گا۔ اگر ہم ایسے مثلث سے شروع کریں جس کے اضلاع \عددی{\kvec{B}}، \عددی{\kvec{C}}، اور \عددی{\kvec{B}+\kvec{C}} ہوں (شکل \حوالہء{A12}) اور  یہ  تین اقدام     لاگو کریں تو  یک بعد دیگرے درج ذیل حاصل ہو گا۔
\begin{enumerate}[1.]
\item
ایک مثلث جس کے اضلاع \عددی{\kvec{B}'}، \عددی{\kvec{C}'}، اور \عددی{(\kvec{B}+\kvec{C})'} ہوں گے جہاں درج ذیل سمتی مساوات مطمئن ہو گی۔
\[\kvec{B}'+\kvec{C}'=(\kvec{B}+\kvec{C})'\]
\item
ایک مثلث جس کے اضلاع \عددی{\kvec{B}''}، \عددی{\kvec{C}''}، اور \عددی{(\kvec{B}+\kvec{C})''} ہوں گے جہاں درج ذیل سمتی مساوات مطمئن ہو گی۔
\[\kvec{B}''+\kvec{C}''=(\kvec{B}+\kvec{C})''\]
(جہاں سمتیات وہی ہیں جو شکل \حوالہء{A11} میں دکھائے گئے ہیں۔)
\item
ایک مثلث جس کے اضلاع \عددی{|\kvec{A}|\kvec{B}''}، \عددی{|\kvec{A}|\kvec{C}''}، اور \عددی{|\kvec{A}|(\kvec{B}+\kvec{C})''} ہوں گے جہاں درج ذیل سمتی مساوات مطمئن ہو گی۔
\begin{align}\label{مساوات_ضمیمہ_جزئیتی_ب}
|\kvec{A}|\kvec{B}''+|\kvec{A}|\kvec{C}''=|\kvec{A}|(\kvec{B}+\kvec{C})''
\end{align}
اب \عددی{|\kvec{A}|\kvec{B}''=\kvec{A}\times\kvec{B}}، \عددی{|\kvec{A}|\kvec{C}''=\kvec{A}\times\kvec{C}}، اور 
\عددی{|\kvec{A}|(\kvec{B}+\kvec{C})''=\kvec{A}\times(\kvec{B}+\kvec{C})} ڈالنے 
سے مساوات \حوالہ{مساوات_ضمیمہ_جزئیتی_ب} درج  ذیل دے گی 
\[\kvec{A}\times\kvec{B}+\kvec{A}\times\kvec{C}=\kvec{A}\times(\kvec{B}+\kvec{C})\]
جو وہ قانون ہے جس کا ثبوت درکار تھا۔
\end{enumerate}
\انتہا{ثبوتءء}
